\chapter{符号和缩写说明}\label{01}

本文的各章节中存在部分通用符号，如无特殊说明，其符号的含义如下表 \ref{symbol} 所示。

\begin{table}[htpb]
\begin{center}
\bicaption{符号及意义}{Symbols and definitions }
\begin{tabular}{l@{\hspace{30pt}}l}
\hline  符号 & 物理意义 \\
\hline
\multicolumn{2}{c}{\textbf{基本数学符号}} \\
\hline
$P^{\top}$ & 矩阵 $P$ 的转置 \\
$P^{-1}$ & 矩阵 $P$ 的逆 \\
$\|x\|$ & 向量 $x$ 的欧几里得范数 \\
$\|G(z)\|_{\infty}$ & 传递函数矩阵 $G(z)$ 的 $H_{\infty}$ 范数 \\
$\mathbb{R}^{n}$ & $n$ 维欧式空间 \\
$\mathbb{R}^{n\times m}$ & $n\times m$ 维实数矩阵 \\
$\operatorname{diag}(\cdot)$ & 以输入元素构成的对角矩阵 \\
$I_n$ & $n$ 阶单位矩阵 \\
$0_{m,n}$ & $m\times n$ 维零矩阵 \\
$\oplus$ & Minkowski 和 \\
$\ominus$ & Pontryagin 差 \\
\hline
\multicolumn{2}{c}{\textbf{系统参数}} \\
\hline
$M$ & 线路上运行的列车总数 \\
$N$ & 沿线车站总数 \\
$H$ & 相邻列车的标称发车间隔（追踪间隔）\\
$S$ & 列车在车站的最小停站时间 \\
$R_{k}$ & 从车站 $k$ 至车站 $k+1$ 的计划运行时间 \\
$c$ / $c_{k}$ & 车站 $k$ 的客流延误率，表示发车间隔对停站时间的影响程度 \\
$a$ / $\beta$ & 单位乘客平均乘降时间影响系数（第3章用 $a$，第5章用 $\beta$）\\
$\delta_{k}$ / $\phi_{k}$ & 车站 $k$ 的平均客流量（第3章用 $\delta_k$，第5章用 $\phi_k$）\\
\hline
\multicolumn{2}{c}{\textbf{列车运行变量}} \\
\hline
$t_{k}^{i}$ & 列车 $i$ 在车站 $k$ 的实际发车时刻 \\
$T_{k}^{i}$ & 列车 $i$ 在车站 $k$ 的计划发车时刻 \\
$r_{k}^{i}$ & 列车 $i$ 从车站 $k$ 至车站 $k+1$ 的实际运行时间 \\
$s_{k}^{i}$ & 列车 $i$ 在车站 $k$ 的实际停站时间 \\
$x_{k}^{i}$ & 列车 $i$ 在车站 $k$ 的实际发车时刻与计划发车时刻的偏差（发车延误），$x_{k}^{i}=t_{k}^{i}-T_{k}^{i}$ \\
$y_{k}^{i}$ & 列车 $i$ 与列车 $i-1$ 在车站 $k$ 的发车间隔相对于标称值 $H$ 的偏差，$y_{k}^{i}=t_{k}^{i}-t_{k}^{i-1}-H$ \\
$u_{k}^{i}$ & 对列车 $i$ 在区间 $k$ 至 $k+1$ 的运行时间施加的重调度调整量 \\
${w_1}_{k}^{i}$ & 列车 $i$ 在区间 $k$ 至 $k+1$ 的运行时间扰动项 \\
${w_2}_{k}^{i}$ & 列车 $i$ 在车站 $k$ 的停站时间扰动项 \\
$w_{k}^{i}$ & 总扰动项，$w_{k}^{i}={w_1}_{k}^{i}+{w_2}_{k+1}^{i}$ \\
$\delta u_{k}^{i}$ & 相邻列车控制量之差，$\delta u_{k}^{i}=u_{k}^{i}-u_{k}^{i-1}$ \\
$\delta w_{k}^{i}$ & 相邻列车扰动之差，$\delta w_{k}^{i}=w_{k}^{i}-w_{k}^{i-1}$ \\
\hline
\multicolumn{2}{c}{\textbf{状态空间模型}} \\
\hline
$X_{k}^{i}$ & $M$ 维状态向量，表征所有列车在各车站的发车延误 \\
$U_{k}^{i}$ & $N$ 维控制向量 \\
$W_{k}^{i}$ & $N$ 维扰动向量 \\
$Y_{k}^{i}$ & 面向规律发车间隔模型的$M$维状态向量 \\
$\overline{U}_{k}^{i}$ & $M$ 维增广控制向量 \\
$\overline{W}_{k}^{i}$ & $M$ 维增广扰动向量 \\
$A$, $A_L$ & 系统状态转移矩阵（$M\times M$）\\
$B$, $B_L$ & 系统控制/扰动输入矩阵（$M\times N$）\\
$D_{1}(a,b)$ & $M\times M$ 变换矩阵，用于描述列车延误演化特性 \\
$D_{2}(a)$ & $M\times N$ 变换矩阵，用于描述控制与扰动输入的特性 \\
$S_{M}$ & $M\times M$ 循环置换矩阵 \\
\hline
\multicolumn{2}{c}{\textbf{控制与优化}} \\
\hline
$Q$, $R$ & 二次型目标函数中的加权矩阵 \\
$p$, $q$ & 加权矩阵中的标量权值参数 \\
$K(i)$ & 状态反馈控制增益矩阵（$N\times M$）\\
$\Gamma$ & $H_{\infty}$ 鲁棒扰动抑制水平 \\
$p$ & 模型预测控制的预测时域 \\
$m$ & 模型预测控制的控制时域 \\
$k$ & 迭代学习控制的批次编号 \\
$\Delta U_{k}$ & 相邻批次控制量之差，$\Delta U_k = U_k - U_{k-1}$ \\
$\Delta W_{k}$ & 相邻批次扰动之差 \\
$\Delta X_{k}$ & 相邻批次状态偏差之差 \\
$b_d$ & 相邻批次间扰动变化的有界上界 \\
\hline
\multicolumn{2}{c}{\textbf{性能指标}} \\
\hline
$\prod_{j}$ & 运行质量指数（OQI），用于衡量时刻表恢复效果 \\
$\sigma_{j}$ & 发车间隔规律性指数，用于衡量发车间隔均一性 \\
$J_1^{k}$ & 名义时刻表恢复控制器的目标函数 \\
$J_2^{k}$ & 规律发车间隔控制器的目标函数 \\
$G$ & 鲁棒MPC中的累计代价函数 \\
\hline
\multicolumn{2}{c}{\textbf{模糊控制系统（第5章）}} \\
\hline
$L$ & 模糊规则总数 \\
$h_l(v)$ & 第 $l$ 条模糊规则的归一化激活强度 \\
$A_l$, $B_l$, $D_l$ & 第 $l$ 个局部模糊子系统的系统矩阵 \\
$M_{ln}$ & 第 $l$ 条规则中第 $n$ 个前提变量的模糊隶属集 \\
\hline
\multicolumn{2}{c}{\textbf{缩写}} \\
\hline
IHSR & 城际高速铁路 (Intercity High-Speed Railway)\\
OQI & 运行质量指数 (Operation Quality Index)\\
FCFS & 先到先服务 (First Come First Served)\\
MIP & 混合整数规划 (Mixed Integer Programming)\\
MPC & 模型预测控制 (Model Predictive Control)\\
RMPC & 鲁棒模型预测控制 (Robust Model Predictive Control)\\
ILC & 迭代学习控制 (Iterative Learning Control)\\
IL-MPC & 迭代学习模型预测控制 (Iterative Learning MPC)\\
ILFPC & 迭代学习模糊预测控制 (Iterative Learning Fuzzy Predictive Control)\\
LMI & 线性矩阵不等式 (Linear Matrix Inequality)\\
T-S & Takagi–Sugeno 模糊模型 \\
PDC & 平行分布补偿 (Parallel Distributed Compensation)\\
\hline
\end{tabular}\label{symbol}
\end{center}
\end{table}

\textbf{符号说明：} 除另有说明，上标 $i$ 表示列车编号（$1 \leq i \leq M$），下标 $k$ 表示车站编号（$1 \leq k \leq N$）。由于城际高铁环形线路的周期性，列车编号与车站编号分别对 $M$ 和 $N$ 取模：列车 $M$ 之后为列车 $1$，车站 $N$ 之后为车站 $1$。第4章采用下标 $j$ 替代 $k$ 表示车站编号，含义相同。
